%==============================================================================
\chapter{الآجال والمُخرَجات}
%==============================================================================

\section{المُخرَجات المنتظرة}
\begin{itemize}
  \item تطبيق وظيفي (شيفرة مصدرية لخادم Java + واجهة أمامية).
  \item خطة اختبارات ونتائج التنفيذ.
  \item ملف تصميم (مخطّطات UML).
  \item تقرير مشروع، ملصق وعرض شفوي.
\end{itemize}

\section{شبكة التقييم (السلّم)}
\begin{xltabular}{\textwidth}{X C{2.5cm}}
\toprule
\textbf{المعيار} & \textbf{النقاط} \\
\midrule
\endhead
سهولة الاستعمال & 4 \\
الوظائف المُختبَرة (خطة الاختبارات + التنفيذ) & 4 \\
\textbf{التقنيات} & \textbf{4} \\
\quad MVC & 0.5 \\
\quad Servlet & 0.5 \\
\quad JSP الواجهة الأمامية & 0.25 \\
\quad ملف XML والتدويل & 0.5 \\
\quad ملف نصّي للإعداد & 0.25 \\
\quad خدمة ويب عميل OAuth & 0.5 \\
\quad خدمة ويب API XML & 0.5 \\
\quad JDBC & 0.5 \\
\quad X.509 / SSL & 0.25 \\
\quad EJB & 0.25 \\
\textbf{التصميم} (UC، أصناف، كائنات، تسلسل، نشاط، حالة، & \\
\quad تفاعل/تعاون، حزمة، مكوّن، نشر) & 4 \\
\textbf{التوثيق} & \textbf{4} \\
\quad التقرير (البنية، الاتّساق، الصيغة) & 2 \\
\quad الملصق & 1 \\
\quad العرض & 1 \\
\midrule
\textbf{المجموع} & \textbf{20} \\
\bottomrule
\end{xltabular}

\section{مخطّطات UML الواجب إنتاجها}
مخطّطات: حالات الاستعمال، الأصناف، الكائنات، التسلسل، النشاط، آلة الحالات،
التفاعل / التعاون، الحزمة، المكوّن والنشر. المحتوى المنتظر لكلٍّ منها مفصّل في
الفصل \ref{chap:conception} المتعلّق بالتصميم.

\section{الرزنامة التقديرية}
تُعبَّر الرزنامة بأسابيع نسبية، مع المُخرَجات الوسيطة المرتبطة بكل مرحلة.
\begin{xltabular}{\textwidth}{>{\bfseries}p{4.5cm} C{2.2cm} X}
\toprule
\textbf{المرحلة} & \textbf{الأسابيع} & \textbf{المُخرَج} \\
\midrule
\endhead
التحليل وكراس الشروط & س1 إلى س2 & كراس شروط مُصادَق عليه \\
التصميم (UML) & س3 إلى س4 & ملف التصميم \\
تطوير نواة CRUD & س5 إلى س7 & النموذج وإدارة الكيانات \\
لوحات القيادة والتنبيهات & س8 إلى س9 & الرزنامة والتنبيهات والخلاصات \\
خدمات الويب والأمان والتدويل & س10 إلى س11 & OAuth، API، X.509 وSSL، XML \\
الاختبار والتحقّق & س12 & خطة اختبارات مُنفَّذة \\
التوثيق والعرض & س13 & التقرير والملصق والمناقشة \\
\bottomrule
\end{xltabular}
